\documentclass[11pt]{article}
\pdfinfoomitdate=1
\pdftrailerid{}
\pdfsuppressptexinfo=15
\usepackage{fullpage,amsmath,amssymb,amsthm}
\usepackage[T1]{fontenc}
\usepackage{lmodern}
\usepackage[title]{appendix}
\usepackage[colorlinks=true,linkcolor=blue,citecolor=blue,urlcolor=blue]{hyperref}
\hypersetup{pdftitle={A counterexample to the E-derivation part of the LNED conjecture},
 pdfauthor={P. Graham},pdfsubject={Mathieu-Zhao subspaces and polynomial differences},
 pdfkeywords={LNED conjecture, E-derivation, Mathieu-Zhao subspace, Bernoulli polynomials, Lean}}
\newtheorem{theorem}{Theorem}
\newtheorem{lemma}[theorem]{Lemma}
\newtheorem{corollary}[theorem]{Corollary}
\newcommand{\Q}{\mathbb{Q}}
\newcommand{\E}{\mathcal{E}}
\title{A counterexample to the $\mathcal{E}$-derivation part of the LNED conjecture}
\author{P. Graham\\ \small\url{https://orcid.org/0009-0007-8493-1505}}
\date{October 25, 2026}
\begin{document}
\maketitle
\begin{abstract}
For a field $K$ of characteristic zero and $c\in K\setminus\{0\}$, let
$\delta_c f(x)=f(x)-f(x+c)$ and $J_c=(x(x-c/2))\subset K[x]$.
We prove that $\delta_c$ is a locally nilpotent $K$-$\E$-derivation
such that the image of $J_c$ under $\delta_c$ contains $x^{2m}$ and excludes $x^{2m+1}$ for every
$m\ge1$, refuting the $\E$-derivation part of Zhao's locally nilpotent
derivation and $\E$-derivation (LNED) conjecture already in one variable.
The proof specializes Zhao's Bernoulli-polynomial membership criterion
for images of quadratic ideals at $x=c/2$.
For any $v$ in an index set $T$, the shift $x_v\mapsto x_v+1$ in
$K[x_i:i\in T]$, fixing the other variables, gives the same
failure of the Mathieu--Zhao condition for the corresponding quadratic ideal.
\end{abstract}

\begin{quotation}\small\noindent
This work has not yet been reviewed by experts.
The proofs, the Lean formalization and the text were produced by AI models
(from OpenAI and Anthropic), working under the author's direction.
\end{quotation}

\noindent 2020 \emph{Mathematics Subject Classification}: 13N15, 11B68.\\
\emph{Keywords}: Mathieu--Zhao subspace, LNED conjecture, Bernoulli polynomials.

\section{Introduction}

A $K$-linear subspace $V$ of a commutative $K$-algebra $A$ is a
\emph{Mathieu--Zhao subspace} if, whenever $w^m\in V$ for all integers
$m\ge1$, every $b\in A$ satisfies $bw^m\in V$ for all sufficiently large
$m$. For example, every ideal has this property. A $K$-$\E$-derivation is a
map $\delta=\operatorname{id}-\phi$, where $\phi$ is a $K$-algebra
endomorphism. It satisfies
\[
 \delta(fg)=\delta(f)g+f\delta(g)-\delta(f)\delta(g).
\]
Such a map is \emph{locally nilpotent} if for each $f$ some positive
iterate $\delta^N(f)$ vanishes.

\begin{theorem}\label{thm:main}
Let $K$ be any field of characteristic zero and $c\in K$ be nonzero.
On $K[x]$, define $\phi_c(f)(x)=f(x+c)$,
$\delta_c=\operatorname{id}-\phi_c$, and $J_c=(x(x-c/2))$.
Then $\delta_c$ is a locally nilpotent $K$-$\E$-derivation, and
\[
 x^{2m}\in\delta_c(J_c),\qquad
 x^{2m+1}\not\in\delta_c(J_c)\qquad(m\ge1).
\]
In particular, $\delta_c(J_c)$ is not a Mathieu--Zhao subspace.
\end{theorem}

Zhao's locally nilpotent derivation and $\E$-derivation (LNED) conjecture,
proposed in \cite[Conjecture~1.6]{zhao-general}\footnote{Statement numbers refer to the arXiv versions:
\cite{essen-zhao-lfed} uses arXiv:1708.05813v2
(2022), \cite{zhao-general} arXiv:1701.05992v2 (2022), and \cite{zhao-univariate}
arXiv:1701.06125v1 (2017).} and restated for
commutative algebras in \cite[Conjecture~1.4]{zhao-univariate},
asserts that every locally
nilpotent $K$-derivation or $K$-$\E$-derivation of a commutative
$K$-algebra, over a field $K$ of characteristic zero, sends every ideal
to a Mathieu--Zhao subspace.
The original formulation also allows noncommutative algebras.
Theorem~\ref{thm:main} disproves the
$\E$-derivation assertion.

Zhao introduced Mathieu subspaces \cite[Definition~1.1]{zhao-origin}
and developed their associative-algebra definition
\cite[Definitions~1.1--1.2]{zhao-assoc}. The terminology Mathieu--Zhao
follows van den Essen \cite[\S\S1--2]{essen-intro}. Van den Essen and Zhao
give a structural criterion for univariate polynomial algebras
\cite[Theorem~4.1]{essen-zhao-univariate}; van den Essen and van Hove
give criteria over an algebraically closed field of characteristic zero
for subspaces containing an ideal of finite codimension
\cite[Theorem~1]{essen-hove}.

In arXiv:1701.06125v1, Zhao proved the ordinary-derivation case for ideals
$(u)$ in $K[x]$
and the translation case when $u=0$, $\deg u\le1$, or $u$ has a
repeated root \cite[Theorem~1.6(1),(3)]{zhao-univariate}.
Zhao's 2017 arXiv version \cite[Theorem~1.6 and the following paragraph]{zhao-univariate}
leaves open nonzero translations with squarefree generators of degree at least two.
The later journal version was not available to us in full; the statement
compared here is arXiv:1701.06125v1, posted in January 2017.
Our generator has the two
distinct roots $0,c/2$.
For quadratic ideals $(x(x-\alpha))$, Zhao gives a Bernoulli-polynomial
membership criterion \cite[Lemma~4.5 and Corollary~4.6]{zhao-univariate};
we give the proof of its monomial instance at $\alpha=c/2$ below.
Zhao's Examples~2.4 and~3.6 \cite{zhao-univariate} already give the even/odd failure of the
Mathieu--Zhao condition for
$x\,d/dx$ and for $\operatorname{id}-\phi$ with $\phi(x)=qx$,
$q\in K^\times$ not a root of unity, respectively; those operators are
not locally nilpotent. The specialization here produces this failure
for the locally nilpotent map $\delta_c$.

Section 2 proves Theorem~\ref{thm:main} and its several-variable
consequence. Section 3 gives the correspondence with the Lean statements,
and Section 4 discusses the remaining questions.

\section{Bernoulli polynomials and the counterexample}


The proof compares the values of a polynomial preimage at the two
roots of $x(x-c/2)$. Zhao's criterion then reduces even- and odd-power
membership to Bernoulli values at one half. Write $B_n(t)$ for the
Bernoulli polynomials with $B_1(t)=t-1/2$, and $B_n=B_n(0)$.
We regard $\Q$ as a subfield of $K$.

\begin{lemma}[Monomial case of Zhao's membership criterion]\label{lem:membership}
For each integer $n\ge0$, put
\[
 H_{n,c}(x)=-\frac{c^n}{n+1}\bigl(B_{n+1}(x/c)-B_{n+1}\bigr).
\]
Then $\delta_c H_{n,c}=x^n$, $H_{n,c}(0)=0$, and
$x^n\in\delta_c(J_c)$ if and only if $H_{n,c}(c/2)=0$.
\end{lemma}
\begin{proof}
The difference identity $B_{n+1}(t+1)-B_{n+1}(t)=(n+1)t^n$
\cite[Eq.~24.4.1]{dlmf} gives the first two statements. If $q$ has degree
$d\ge1$ and leading coefficient $q_d$, then $\delta_c q$ has leading
term $-dc q_d x^{d-1}$. Consequently the kernel of $\delta_c$ consists
of constants, and every preimage of $x^n$ is $H_{n,c}+k$ for some $k\in K$.
Since $c/2\ne0$, membership in $J_c$ is equivalent to vanishing at $0$
and $c/2$.
The first root forces $k=0$, and the second gives the equivalence.
This is the monomial instance of Zhao's criterion
\cite[Lemma~4.5 and Corollary~4.6]{zhao-univariate}.
\end{proof}

In particular, the value at the second root is
\begin{equation}\label{eq:half}
 H_{n,c}(c/2)=-\frac{c^n}{n+1}(2^{-n}-2)B_{n+1},
\end{equation}
by $B_j(1/2)=(2^{1-j}-1)B_j$ \cite[Eq.~24.4.27]{dlmf}.

\begin{proof}[Proof of Theorem~\ref{thm:main}]
Since $\phi_c$ is multiplicative,
$\delta_c(fg)=fg-\phi_c(f)\phi_c(g)$, which expands to the
product identity in Section 1.
The degree computation in Lemma~\ref{lem:membership} shows that
$\delta_c^{d+1}q=0$ whenever $\deg q\le d$, proving local nilpotence.

The standard Bernoulli facts give $B_{2m+1}=0$ for $m\ge1$ and
$B_{2k}\ne0$ for $k\ge1$ \cite[Eq.~24.2.2]{dlmf}.
For $n=2m$ with $m\ge1$, equation~\eqref{eq:half} vanishes since
$B_{2m+1}=0$. For $n=2m+1$, it is nonzero since $B_{2m+2}\ne0$ and
$2^{-(2m+1)}-2\ne0$. Lemma~\ref{lem:membership} gives the two power
assertions. Taking $w=x^2$ and $b=x$, one has $bw^m=x^{2m+1}$ outside
the image for every $m\ge1$, violating the Mathieu--Zhao condition.
\end{proof}

\begin{corollary}\label{cor:variables}
Let $K$ have characteristic zero, let $T$ be any index set, and choose
$v\in T$. In $K[x_i:i\in T]$, let $\phi$ send $x_v$ to $x_v+1$ and fix
the other variables. Then $\delta=\operatorname{id}-\phi$ is a locally
nilpotent $K$-$\E$-derivation, and the image of $(x_v(x_v-1/2))$ contains $x_v^{2m}$ and
excludes $x_v^{2m+1}$ for every $m\ge1$. The image is not a Mathieu--Zhao subspace.
\end{corollary}
\begin{proof}
Regarding each polynomial as a polynomial in $x_v$, its degree in that
variable drops under $\delta$. The univariate even-power preimages embed
in this algebra. If an odd power had a preimage in the indicated ideal,
evaluation of all other variables at zero, which commutes with $\delta$,
would give a preimage in $(x_v(x_v-1/2))\subset K[x_v]$, contrary to
Theorem~\ref{thm:main} with $c=1$. Take $w=x_v^2$ and $b=x_v$.
\end{proof}

\section{Formalization}

The Lean formalization proves Theorem~\ref{thm:main} and
Corollary~\ref{cor:variables}, using the two-sided form $bw^me\in V$
of the Mathieu--Zhao definition. For $V=\delta_c(J_c)$ this condition is
\[
 \forall w,\ (\forall m\ge1,\ w^m\in V)\ \Longrightarrow\
 \forall b,e,\ \exists N,\ \forall m\ge N,\ bw^me\in V,
\]
where $w,b,e\in K[x]$ and $m,N\in\mathbb N$; the corollary uses the same
condition in $K[x_i:i\in T]$. The formalization verifies linearity, local
nilpotence, the $\E$-derivation identity and both power assertions, and
proves the negation of this condition. Its Bernoulli value and
nonvanishing lemmas use Mathlib's Bernoulli and zeta-value results.
For commutative algebras this condition is equivalent to that in Section 1:
one direction takes $e=1$, and the other absorbs $e$ into $b$.
The equivalence follows from commutativity and is not part of the formalization.
The Lean sources and checking notes, including the exact declarations,
specialized membership criterion and Bernoulli identities, are available as
\href{https://github.com/6/clank-lab/tree/lned-counterexample-v1/results/lned-counterexample/lean}{supplementary material}.

\section{Concluding remarks}

For a fixed nonzero $c$, the natural remaining problem is to classify
squarefree polynomials $u$ for which $\delta_c((u))$ is a
Mathieu--Zhao subspace. Theorem~\ref{thm:main} treats $u=x(x-c/2)$.
The Mathieu--Zhao property depends on $u$: for example,
$u=x^2$ gives a Mathieu--Zhao image by Zhao's repeated-root case
\cite[Theorem~1.6(3)]{zhao-univariate}, whereas our squarefree quadratic does not.
For $c=0$, the image of every ideal under $\delta_c$ is zero.

The locally finite derivation and $\E$-derivation (LFED) conjecture \cite[Conjecture~1.5]{zhao-general} concerns the image of the
whole algebra under a map whose iterates of each element span a
finite-dimensional subspace. Here $\delta_c(K[x])=K[x]$, since
Lemma~\ref{lem:membership} gives preimages of all monomials; hence this
example satisfies that assertion. In the same arXiv version, Zhao proves the LFED conjecture for $K[x]$
\cite[Theorem~1.5]{zhao-univariate}, and both conjectures for
finite-dimensional characteristic-zero algebras
\cite[Proposition~2.4]{zhao-algebraic}. Van den Essen and Zhao prove LFED
for rational-function and formal-power-series algebras
\cite[Corollaries~2.5 and~2.8]{essen-zhao-lfed}, and for Laurent
formal-power-series algebras \cite[Theorems~3.1--3.2]{essen-zhao-lfed}. The ordinary-derivation
part of LNED is unaffected by our counterexample.

\begin{thebibliography}{10}
\raggedright
\bibitem{dlmf}
F. W. J. Olver et al. (eds.),
\href{https://dlmf.nist.gov/}{\emph{NIST Digital Library of Mathematical
Functions}}, Release 1.2.8, September 15, 2026.

\bibitem{essen-intro}
Arno van den Essen, \emph{An introduction to Mathieu subspaces},
\href{https://arxiv.org/abs/1907.06107v1}{arXiv:1907.06107v1} (2019).

\bibitem{essen-hove}
Arno van den Essen and Loes van Hove, \emph{Mathieu-Zhao spaces of
polynomial rings},
\href{https://arxiv.org/abs/1907.06106v1}{arXiv:1907.06106v1} (2019).

\bibitem{essen-zhao-univariate}
Arno van den Essen and Wenhua Zhao, \emph{Mathieu subspaces of univariate
polynomial algebras}, \emph{Journal of Pure and Applied Algebra}
\textbf{217} (2013), no.~7, 1316--1324,
\href{https://doi.org/10.1016/j.jpaa.2012.10.006}{doi:10.1016/j.jpaa.2012.10.006}.
\href{https://arxiv.org/abs/1012.2017v1}{arXiv:1012.2017v1}.

\bibitem{essen-zhao-lfed}
Arno van den Essen and Wenhua Zhao, \emph{On images of locally finite
derivations and $\mathcal E$-derivations}, \emph{Journal of Pure and
Applied Algebra} \textbf{223} (2019), no.~4, 1689--1698,
\href{https://doi.org/10.1016/j.jpaa.2018.07.002}{doi:10.1016/j.jpaa.2018.07.002}.
\href{https://arxiv.org/abs/1708.05813v2}{arXiv:1708.05813v2} (2022).

\bibitem{zhao-origin}
Wenhua Zhao, \emph{Generalizations of the image conjecture and the
Mathieu conjecture}, \emph{Journal of Pure and Applied Algebra}
\textbf{214} (2010), no.~7, 1200--1216,
\href{https://doi.org/10.1016/j.jpaa.2009.10.007}{doi:10.1016/j.jpaa.2009.10.007}.
\href{https://arxiv.org/abs/0902.0212v3}{arXiv:0902.0212v3}.

\bibitem{zhao-assoc}
Wenhua Zhao, \emph{Mathieu subspaces of associative algebras},
\emph{Journal of Algebra} \textbf{350} (2012), 245--272,
\href{https://doi.org/10.1016/j.jalgebra.2011.09.036}{doi:10.1016/j.jalgebra.2011.09.036}.
\href{https://arxiv.org/abs/1005.4260v2}{arXiv:1005.4260v2}.

\bibitem{zhao-algebraic}
Wenhua Zhao, \emph{The LNED and LFED conjectures for algebraic algebras},
\emph{Linear Algebra and its Applications} \textbf{534} (2017), 181--194,
\href{https://doi.org/10.1016/j.laa.2017.08.016}{doi:10.1016/j.laa.2017.08.016}.
\href{https://arxiv.org/abs/1701.05990v1}{arXiv:1701.05990v1}.

\bibitem{zhao-general}
Wenhua Zhao, \emph{Some open problems on locally finite or locally
nilpotent derivations and $\mathcal E$-derivations},
\emph{Communications in Contemporary Mathematics} \textbf{20} (2018),
no.~4, 1750056, 25 pp.,
\href{https://doi.org/10.1142/S0219199717500560}{doi:10.1142/S0219199717500560}.
\href{https://arxiv.org/abs/1701.05992v2}{arXiv:1701.05992v2} (2022).

\bibitem{zhao-univariate}
Wenhua Zhao, \emph{Images of ideals under derivations and
$\mathcal E$-derivations of univariate polynomial algebras over a field of
characteristic zero},
\href{https://arxiv.org/abs/1701.06125v1}{arXiv:1701.06125v1} (2017).
Published version: \emph{Journal of Algebra and Its Applications}
\textbf{23} (2024), no.~6, 2450128,
\href{https://doi.org/10.1142/S0219498824501287}{doi:10.1142/S0219498824501287}.
\end{thebibliography}
\end{document}
